\chapter{Conclusion} \label{Chap8}

This project sets out to develop and evaluate
machine learning approaches for CH\textsubscript{4} gap-filling, forecasting, and scenario
projection at the North Wyke Farm Platform. All three have been delivered: a benchmarked
gap-filling pipeline (Chapter~\ref{Chap4}) improves combined hourly $R^2$ from 0.055 (MDS) to
0.515 (TabICLv2); a forecasting benchmark (Chapter~\ref{Chap5}), the first applied to this
ecosystem-scale signal at a managed temperate grassland, beats a seasonal-mean baseline by a wide
margin (MASE$=$0.6908, TabPFN); and a scenario-projection architecture (Chapter~\ref{Chap6})
produces distinguishable CH\textsubscript{4} responses to livestock, grazing, and climate
perturbations under CMIP6-derived trajectories to 2050.

Across the forecasting and scenario analyses, livestock-associated variables
emerged as the strongest recurring influence on predicted
CH\textsubscript{4} flux. Permutation importance placed livestock units, total
liveweight, and cattle density among the leading predictors at T4 and T9
(Chapter~\ref{Chap5}), while increasing livestock density produced the largest
scenario response: 100--184\% at each tower's historical maximum, compared
with at most 20\% for extended grazing and 8\% for fertiliser
(Chapter~\ref{Chap6}). Together, these results demonstrate a consistent
livestock-related signal, although they do not establish causality.

Tabular foundation models are the second recurring finding: TabICLv2 leads gap-filling and TabPFN
leads forecasting, both pretrained via in-context learning rather than fitting parameters to
NWFP's own severely incomplete record, a structural advantage as discussed in
Chapter~\ref{Chap7}.

It is worth being precise about what has, and has not, been shown. Gap-filling's improvement over
MDS is real but its absolute fit remains modest ($R^2=0.515$), and weakest specifically at the
livestock-dominant catchments this dissertation's central finding is about (Chapter~\ref{Chap7}).
Forecasting reliably beats a seasonal-mean baseline but not during CH\textsubscript{4}'s dominant
spike events (MASE 3.285 vs.\ 0.524 on ordinary days). Scenario projection produces
distinguishable, interpretable projections while carrying the limitation that no future scenario
can be checked against ground truth. 

Taken together, this dissertation delivers a benchmarked,
uncertainty-quantified foundation for CH\textsubscript{4} forecasting and scenario projection at a
managed UK grassland site that did not previously exist in the literature.
